
%\needspace{4\baselineskip}
\color{uniblue}{\subsection{Фиксация положения выключателей}}
\color{black}

\begin{enumerate}[label=\arabic{section}.\arabic{subsection}.\arabic{enumi}, labelsep=4pt, leftmargin=0pt, itemindent=57pt, itemsep=0pt, parsep=0pt] % если что можно крутить itemsep

\item
Функциональный блок фиксации положения выключателей (ФПВ) предназначен для определения состояния выключателей на основании технологической информации:

\begin{itemize}
    \item состояния блок-контактов выключателя;
    \item положения разъединителей в цепи выключателя;
    \item информации о выводе в ремонт выключателя.
\end{itemize}

\item
Контроль положения разъединителей при формировании сигналов включенного или отключенного положения выключателей (например, \textit{«В1 включен»} и \textit{«В2 отключен»} для выключателя В1) вводится программным ключом \textbf{«Контр. Р В1»}. Также имеется возможность выбрать количество разъединителей, участвующих в формировании сигналов включенного или отключенного положения выключателей, для выключателя В1 она задается программным переключателем \textbf{«Схема Р В1»}.

\item
Также функция ФПВ пофазно контролирует состояние блок-контактов выключателя и в случае расхождения полюсов формирует сигнал \textit{«Внутр. пуск. ЗНФ В1»}.

\item
Функциональная схема логики ФПВ В1 приведена на рисунке \ref{pict:fpv}. Функциональная схема логики ФПВ В2 выполняется аналогично схеме логики ФПВ В1.

\vspace{2mm}
\begin{figure}[H]
\centering
\includegraphics[width=0.9\textwidth,height=0.9\textheight,keepaspectratio]{\fbpath/761.3381 ФПВ/img/fpv.pdf}
\captionof{figure}{Функциональная схема логики ФПВ В1}\label{pict:fpv}
\end{figure}

\end{enumerate}